\chapter{Discussion}
\label{ch:discussion}

This chapter interprets the results of \Cref{ch:results}, against both the original question of the thesis: whether terrain can be recognised from intrinsic sensors alone and the literature in \Cref{ch:background}. \Cref{sec:disc_performance} reads the headline numbers and the residual error structure; \Cref{sec:disc_generalisation} covers cross-run generalisation and the robustness diagnostics; \Cref{sec:disc_deployment} translates the per-class results into operational implications; and \Cref{sec:disc_limitations} sets out what the experiments do and do not support.

\section{Recognition Performance and Terrain Separability}
\label{sec:disc_performance}

The headline farm result, macro-$F_1 \approx 0.86$--$0.87$ under LORO (\Cref{tab:farm_master}), is the main empirical answer to the thesis question, and it comes with a clear structure underneath it.

The first thing the model comparison shows is that, at this dataset scale, the classifier family barely matters. All six families sit in a $0.84$--$0.87$ band that is narrower than the within-family cross-fold standard deviation, the paired-Wilcoxon test returns no significant pair (\Cref{fig:significance_farm}), and the 1D-CNN trained directly on the raw $5\times 200$ window matches the engineered-feature baselines within one standard deviation on every dataset. Most of the discriminative information lives in a low-dimensional projection of the vibration spectrum and the cross-sensor correlations, and any reasonable classifier picks it up. This matches the recurring finding in the IMU classification literature that careful features keep classical learners competitive with deep models at moderate dataset sizes~\cite{Weiss2011,Otsu2016,Fawaz2019}.

What does vary is the per-class structure. \texttt{Dirt\_track} is essentially solved on every fold ($F_1 = 0.98$), \texttt{soil} reaches $0.86$, and \texttt{grass} sits at $0.76$ with the highest cross-fold variance (\Cref{tab:farm_perclass}); the SVM confusion matrix puts almost all of the \texttt{grass} error mass on \texttt{soil} (\Cref{fig:farm_cm}). This makes sense as a compacted track produces a high-amplitude, broadband signature the other two classes do not, while short grass on firm soil and bare loose soil overlap on the same low-amplitude, low-frequency portion of the spectrum and excite the chassis through similar pitch, roll and vertical-bounce modes. The cross-sensor correlations identified by MI (\Cref{sec:feat_mi_pruning}) are what separate them in practice, but the margin is intrinsically narrow.

The campus comparison helps locate where this difficulty comes from. The same pipeline reaches $F_1 \ge 0.98$ on the three-class campus subset (\Cref{tab:campusB_master}), and \Cref{fig:per_class_farm_vs_campus} shows that the $\sim 0.20$ farm-minus-campus gap on \texttt{grass} holds across every classical family. The campus three-class problem spans the full roughness range: smooth asphalt, vegetated lawn, compacted dirt, with no analogue of the \texttt{grass}--\texttt{soil} pair; the farm data is just intrinsically closer in vibration space. Lifting the farm headline therefore looks more like a class-geometry problem than a model-capacity one.

Speed had less impact initially expected. The campus A--to--B comparison adds $0.003$--$0.026$ macro-$F_1$ from the low-constant-speed regime on top of the three-class label set (\Cref{tab:campus_A3_vs_B3}), an order of magnitude less than the class-imbalance step, and the speed-only feature ablation (\Cref{tab:speed_only}) reaches only $F_1 \approx 0.49$--$0.50$, well below the full pipeline. Removing the speed columns while keeping the vibration features costs essentially nothing ($\le 0.003$ in either direction). The classifier registers speed but does not lean on it. The velocity--vibration confound flagged in \Cref{sec:bg-vel} is present and points in the expected direction, but it is not what drives the headline at the operating point sampled here.

\section{Cross-Run Generalisation and Robustness}
\label{sec:disc_generalisation}

LORO (\Cref{sec:eval_protocol}) prevents within-run auto-correlation from leaking into the test fold, so the numbers in \Cref{ch:results} are real cross-run estimates. But the protocol measures one specific generalisation axis, and the perturbation experiments add further nuance.

The axis LORO measures is cross-day, same-site. The five farm runs covered the same Taastrup field on two adjacent days, and the four campus sessions covered the same campus route on different days. Cross-field, cross-platform and cross-weather generalisation are not tested. The clearest place this matters is the campus 5-class collapse on \texttt{cobblestone} and \texttt{muddy\_dirt\_track} ($F_1 = 0.31$ and $0.43$): both classes are recorded almost exclusively on a single session (\Cref{tab:fold_quality_campus}), so when that session is held out, the training fold contains barely any material for them. Speed conditioning compounds it, with per-class train--test TVDs of $0.56$--$0.86$. The takeaway is a methodological one: LORO is only as good as the diversity of the underlying campaign, and the farm campaign was deliberately designed to put every class into every run for exactly this reason (\Cref{sec:farm_dataset}).

The sensor dropout study (\Cref{fig:dropout_compare}) shows a sharp split between transient and permanent failure. Zero-ing the accelerometer or gyroscope columns on the test fold only costs the farm classifier $\Delta F_1 \approx -0.3$, but removing the same channels permanently and retraining costs only $-0.01$ to $-0.02$. The model can express the terrain signature in either inertial group on its own, since the pitch and roll rates and vertical accelerometer pick up the same physical excitation. But if it has been trained on both and then loses one without notice, the standardised zeros land far outside the in-distribution feature space and the decision boundary breaks. For deployment, that distinction matters: graceful degradation needs either retraining when a sensor fails or an inference-time imputation that keeps the input on the trained manifold; a silent zero is the worst case. Odometry is free in both modes, which lines up with the speed-only ablation above.

Noise sensitivity (\Cref{fig:noise_sensitivity}) goes the opposite way from what one might expect. The farm classifier degrades from $0.84$ to $0.59$ as $\sigma$ goes from $0$ to $2$; the near-saturated Dataset~B model starts higher ($0.86$) but degrades faster (to $0.47$). Dataset~B's clean accuracy rests on a small number of sharp feature contrasts between physically distinct surfaces, and noise erodes them quickly; the farm classifier spreads its weight across many cross-sensor correlations between similar classes, and the collective signal survives noise better. Higher clean accuracy does not automatically mean higher robustness.

The smaller-effect findings are worth flagging briefly. Tuned baselines stay within $\pm 0.01$ macro-$F_1$ of the untuned ones (\Cref{sec:res_models}), and top-MI versus all-features swings by $\sim 0.03$ either way. The CUSUM detector (\Cref{tab:transition_cusum}) gives the explicit, calibratable false-alarm rate that the implicit $n_{\text{stable}}$ rule cannot: for the slower classical models (SVM, LR, RF) it cuts $\sim 2$--$9\,\text{s}$ of mean delay at the cost of roughly $10$ percentage points of detection rate, while for the already-fast MLP and CNN the implicit rule is essentially as good. The operating curve in \Cref{fig:cusum_operating} matches the Wald approximation~\cite[Eq.~6.74]{Blanke2003}, false-alarm rate dropping roughly exponentially in $h$, mean delay roughly linear, so the threshold can be calibrated to a target budget rather than re-tuned per deployment.

\section{Operational Implications for Field Deployment}
\label{sec:disc_deployment}

An intercropping robot spends most of its time on \texttt{grass} (between rows) and \texttt{soil} (within the plot); it only crosses \texttt{dirt\_track} when entering, leaving, or returning to charge. The per-class farm numbers should be read with that mix in mind rather than as three equal entries in the macro average.

That puts the binding number not at the $0.87$ headline but at the \texttt{grass}--\texttt{soil} margin. \texttt{Dirt\_track} at $F_1 = 0.98$ is more than sufficient for the role it plays, a reliable transit-versus-field flag, but the in-field discrimination is between \texttt{soil} ($0.86$) and \texttt{grass} ($0.76$). Pushing those specifically up, by adding terrain-specific signals such as per-wheel tyre-slip or motor-current variance, by integrating posteriors over several windows, or by widening the in-field training distribution, would have a larger operational effect than lifting the headline by switching to a heavier model on the same features.

The structure of the residual errors is also worth noting. \Cref{fig:farm_cm} puts almost all of the \texttt{grass} confusion on \texttt{soil} rather than on \texttt{dirt\_track}: when the classifier is wrong, it tends to be wrong between the two physically similar in-field surfaces rather than between an in-field surface and the transit track. Whether this is operationally easier to absorb than a uniform error pattern depends on the downstream controller, which is outside the scope of this thesis; the limited statement that the present data supports is that the error mass aligns with the physical similarity structure of the classes, consistent with the interpretation in \Cref{sec:disc_performance}.

The transition detector is fast enough to be useful at agricultural cruising speeds. The CUSUM operating point at one false alarm per minute (\Cref{tab:transition_cusum}) gives $3$--$6\,\text{s}$ mean detection delay. At the operating speed of $\approx\SI{0.5}{\meter\per\second}$ that is $1.5$--$3\,\text{m}$ of travel between the physical transition and the alarm, roughly one robot footprint and well within the latency budget of a slow weeding manoeuvre.

Finally, the sensor set the pipeline needs is exactly the one Ricbot already carries (\Cref{ch:setup}), and it requires no GNSS, no camera, no LiDAR at inference time, and the pipeline is invariant to the lighting, dust and glare that defeat exteroceptive sensors in field conditions~\cite{Reina2016}. Within the limits set out next, that is the lightweight, always-available channel the thesis set out to demonstrate.

\section{Limitations}
\label{sec:disc_limitations}

The findings above need to be qualified by the scope of the experiments that produced them.

\paragraph{Single platform, single field, single weather window.} All measurements were taken with one robot (Ricbot, $25$--$30\,\text{kg}$, differential drive, \Cref{ch:setup}), at one Taastrup field, across two adjacent dry sunny April days (\Cref{tab:farm_sessions}). A heavier or differently suspended chassis would change the vibration signature, and a different soil type, grass species or wet-soil regime would shift the within-class distributions. The reported numbers generalise across days and operator-driven speed variation within the sampled site; extrapolation to other sites, weather, or platforms is not supported by the present data.

\paragraph{Scope of generalisation testing.} LORO holds the physical site fixed and varies only the session, so the reported numbers are cross-day, same-site estimates. Three natural extensions sit outside this work. Routes were permuted between farm runs to broaden the within-class distribution, but this is not Leave-One-Field-Out; a true cross-field test would need a second farm campaign. Additional farm sessions covering different vegetation stages or weather windows would also have been valuable, but the campaign slipped into late April (\Cref{sec:project-plan}) and further recordings could not be fitted into the remaining timeline. Cross-domain transfer from campus to farm is not attempted either, because the label spaces only partially overlap (\Cref{sec:terrain_classes_and_labels}); the campus-versus-farm shared-class diagnostic in \Cref{sec:res_models} is a per-class substitute, not a transfer test. All farm results therefore characterise a single Taastrup field under dry, sunny conditions in a single fortnight.

\paragraph{Narrow speed range.} The farm campaign was driven at the low, near-constant speed natural for intercropping (\Cref{sec:bg-vel}). The campus A--to--B comparison and the speed-only ablation jointly argue that the velocity confound is small at the sampled operating point (\Cref{sec:disc_performance}), but neither addresses higher cruising speeds or transitions taken at speed.

\paragraph{Dataset size.} The farm and campus datasets contain $\sim 1.6\text{k}$ and $\sim 2.3\text{k}$ windows (\Cref{tab:window_counts}). At this scale the engineered features carry most of the discriminative weight, and the 1D-CNN consistently shows a larger worst-fold spread than the classical baselines (\Cref{sec:res_models}), consistent with operating at the lower edge of the regime in which end-to-end deep models start to pay off~\cite{Wang2017,Fawaz2019}. Whether a much larger dataset would shift the ranking is not testable here.

\paragraph{Transition detection evaluated only on the farm.} CUSUM and its operating curve (\Cref{sec:res_transition,fig:cusum_operating}) are reported only on the farm. The farm is the deployment target, so it was the priority for the explicit detector; CUSUM itself was also added late in the project(\Cref{sec:project-plan}), and a campus rerun was outside the remaining timeline. The campus $n_{\text{stable}}$ baseline does show a heavy-tailed delay driven by a single uncertain transition on the snowy 23~Feb log (\Cref{tab:transition_nstable}), exactly the failure mode an explicit threshold is meant to control, so extending CUSUM to the campus dataset is the natural next step.

\paragraph{Conservative labelling.} The labelling procedure in \Cref{sec:terrain_classes_and_labels} discards ambiguous boundary regions in exchange for cleaner per-class assignments, so gradual transitions (slow firm-to-loose soil changes, for example) are excluded rather than modelled. The transition analysis treats each change as a point event; modelling boundaries as continuous mixture regions would need richer ground-truth annotation.

\paragraph{Single chassis-mounted IMU.} The pipeline uses one MEMS IMU on the central chassis. Per-wheel suspension sensors, motor-current measurements, microphones, or several distributed IMUs all appear in the broader literature~\cite{DuPont2008,Ojeda2006,Bermudez2012} and could plausibly add information cheaply, especially for the \texttt{grass}--\texttt{soil} pair where tyre slip and motor load differ even when chassis vibration overlaps. This is left as future work.

Within these limits, the experiments of \Cref{ch:results} give a self-consistent and reproducible picture of what an intrinsic-only terrain-classification pipeline can do on a small wheeled agricultural robot, evaluated on the harder of the two campaigns at the operating point for which the pipeline was designed.